% Wersja źródłowa: audyt II, jednostka 16
\setcounter{section}{15}
\section{Funkcjonały wypukłe i podstawy rachunku wariacyjnego}
\label{sec:wariacyjny16}

\subsection{Od pochodnej Frécheta do problemu minimum}

W poprzednich rozdziałach badaliśmy lokalną strukturę funkcji
różniczkowalnych na przestrzeniach Banacha. Rachunek wariacyjny pyta
o istnienie i charakter minimizatorów funkcjonału $J:E\to\R$,
często na nieskończenie wymiarowej przestrzeni funkcji.
W tym rozdziale $H$ oznacza rzeczywistą przestrzeń Hilberta,
a $E$ rzeczywistą przestrzeń Banacha. Pisząc $\langle\cdot,\cdot\rangle$
w $H$, przyjmujemy liniowość w pierwszym argumencie.

\begin{definicja}[Wypukłość funkcjonału]
\label{def:convex-fun16}
Niech $E$ będzie rzeczywistą przestrzenią liniową,
a $C\subseteq E$ zbiorem wypukłym.
Funkcjonał $J:C\to(-\infty,+\infty]$ nazywamy
\emph{wypukłym}, jeżeli
$J((1-t)x+ty)\le(1-t)J(x)+tJ(y)$ dla wszystkich $x,y\in C$
i $0<t<1$. Dla $t=0,1$ przyjmujemy równość, bez
mnożenia zera przez nieskończoność.
Funkcjonał nazywamy \emph{właściwym}, gdy jego dziedzina
efektywna $\operatorname{dom}J=\{x\in C:J(x)<+\infty\}$
jest niepusta. Z wypukłości $J$ wynika wypukłość
$\operatorname{dom}J$, ponieważ kombinacja wypukła
dwóch skończonych wartości jest skończona.
Funkcjonał wypukły jest \emph{ściśle wypukły}, jeśli nierówność
jest ostra dla różnych $x,y\in\operatorname{dom}J$
i $0<t<1$. W przypadku funkcjonału rzeczywistego
$\operatorname{dom}J=C$.
\end{definicja}

\begin{twierdzenie}[Nierówność styczna]
\label{thm:convex-derivative16}
Niech $E$ będzie rzeczywistą przestrzenią Banacha,
$U\subseteq E$ niepustym otwartym zbiorem wypukłym, a
$J:U\to\R$ funkcjonałem różniczkowalnym w sensie
Frécheta w każdym punkcie $U$.
Następujące warunki są równoważne:
\begin{enumerate}
\item $J$ jest wypukły na $U$;
\item $J(y)\ge J(x)+DJ(x)(y-x)$ dla wszystkich $x,y\in U$.
\end{enumerate}
Jeżeli $DJ(x_0)=0$, gdzie $x_0\in U$, a $J$ jest wypukły,
to $x_0$ jest globalnym minimizatorem $J$ na $U$.
Odwrotnie, jeśli $x_0$ jest lokalnym minimizatorem
różniczkowalnego $J$ i $x_0\in U$, to $DJ(x_0)=0$.
\end{twierdzenie}
\begin{proof}
Dla ustalonych $x,y\in U$ niech
$\phi(t)=J(x+t(y-x))$, $0\le t\le1$.
Jeżeli $J$ jest wypukły, to dla $0<t\le1$
z definicji wypukłości
$[\phi(t)-\phi(0)]/t\le\phi(1)-\phi(0)$.
Z reguły łańcuchowej (Twierdzenie~\ref{thm:chain14})
pochodna $\phi$ w zerze wynosi $DJ(x)(y-x)$
(otwartość $U$ pozwala przedłużyć parametr nieco poza zero).
Przechodząc z $t\downarrow0$, otrzymujemy
$DJ(x)(y-x)\le J(y)-J(x)$.

Odwrotnie, ustalmy $x,y\in U$, $0<t<1$ i
$z=(1-t)x+ty\in U$.
Nierówność z punktu 2 daje
$J(x)\ge J(z)+DJ(z)(x-z)$ oraz
$J(y)\ge J(z)+DJ(z)(y-z)$.
Mnożąc je odpowiednio przez $1-t$ i $t$ oraz dodając,
otrzymujemy
\[
 (1-t)J(x)+tJ(y)\ge J(z)+DJ(z)((1-t)x+ty-z)=J(z).
\]
Dla $t=0,1$ zachodzi równość, więc $J$ jest wypukły.
Jeżeli $DJ(x_0)=0$, nierówność styczna zapewnia
$J(y)\ge J(x_0)$. Jeśli zaś $x_0$ jest lokalnym minimum,
dla każdego $h\in E$ i dostatecznie małych $|t|$
punkt $x_0+th$ należy do $U$ i do otoczenia lokalnego
minimum. Funkcja $t\mapsto J(x_0+th)$ ma zatem
minimum w zerze, a jej pochodna w zerze, równa
$DJ(x_0)h$, musi znikać. Ponieważ dotyczy to
każdego $h\in E$, mamy $DJ(x_0)=0$.
\end{proof}

\begin{definicja}[Silna wypukłość]
\label{def:strong-convex16}
Niech $H$ będzie rzeczywistą przestrzenią Hilberta
i $\mu>0$. Funkcjonał $J:H\to\R$ jest
$\mu$-silnie wypukły, jeśli $x\mapsto J(x)-\frac\mu2\|x\|^2$
jest wypukły.
\end{definicja}

\begin{stwierdzenie}[Koercywność z silnej wypukłości]
\label{prop:strong-coercive16}
Niech $H$ będzie rzeczywistą przestrzenią Hilberta,
$\mu>0$, a $J:H\to\R$ funkcjonałem $\mu$-silnie wypukłym
i różniczkowalnym w sensie Frécheta na $H$. Wtedy dla $x,y\in H$
zachodzi
\[
J(y)\ge J(x)+DJ(x)(y-x)+\frac\mu2\|y-x\|^2.
\]
W szczególności $J$ jest koercywny, tj.
$J(y)\to+\infty$, gdy $\|y\|\to\infty$,
i ma co najwyżej jeden minimizator.
\end{stwierdzenie}
\begin{proof}
Stosujemy Twierdzenie~\ref{thm:convex-derivative16}
do funkcjonału $G(z)=J(z)-\mu\|z\|^2/2$.
Z rozwinięcia
$\|x+h\|^2=\|x\|^2+2\langle x,h\rangle+\|h\|^2$
wynika $DG(x)h=DJ(x)h-\mu\langle x,h\rangle$.
Nierówność styczna dla $G$ daje
\[
 J(y)-\tfrac\mu2\|y\|^2
 \ge J(x)-\tfrac\mu2\|x\|^2
       +DJ(x)(y-x)-\mu\langle x,y-x\rangle .
\]
Po przeniesieniu składników kwadratowych na prawą stronę
i wykorzystaniu tożsamości
$\|y\|^2-\|x\|^2-2\langle x,y-x\rangle=\|y-x\|^2$
otrzymujemy żądaną nierówność.
Dla $x=0$ mamy
$J(y)\ge J(0)-\|DJ(0)\|\|y\|+\mu\|y\|^2/2$,
co zapewnia koercywność. Jeżeli $x\ne y$ są minimizatorami,
to wypukłość $G$ przy $t=1/2$ i tożsamość
$\tfrac12(\|x\|^2+\|y\|^2)-\|(x+y)/2\|^2
=\|x-y\|^2/4$ dają
$J((x+y)/2)\le [J(x)+J(y)]/2-\mu\|x-y\|^2/8$,
sprzecznie z minimalnością.
\end{proof}

\begin{przyklad}[Niezbędne jednostajne oszacowanie hesjanu]
\label{ex:variational-hessian16}
Niech $H$ będzie rzeczywistą przestrzenią Hilberta
i $J\in C^2(H,\R)$. Jeżeli istnieje $\mu>0$ takie, że
$D^2J(x)[h,h]\ge\mu\|h\|^2$ dla wszystkich $x,h\in H$,
to $J$ jest $\mu$-silnie wypukły.
Istotnie, dla $x,y\in H$ i $h=y-x$ funkcja
$\phi(t)=J(x+th)-\frac\mu2\|x+th\|^2$
ma drugą pochodną
$\phi''(t)=D^2J(x+th)[h,h]-\mu\|h\|^2\ge0$.
Zatem $\phi'$ jest niemalejąca. Dla $0<t<1$
z twierdzenia o wartości średniej na $[0,t]$ i $[t,1]$
wynika
$(\phi(t)-\phi(0))/t\le(\phi(1)-\phi(t))/(1-t)$.
Po pomnożeniu przez $t(1-t)$ i uporządkowaniu
otrzymujemy $\phi(t)\le(1-t)\phi(0)+t\phi(1)$,
czyli wypukłość funkcjonału $J-\mu\|\cdot\|^2/2$.
Sama dodatnia określoność $D^2J(0)$, nawet wraz
z $DJ(0)=0$, nie gwarantuje lokalnego minimum:
\hyperref[ex:positive-hessian15]{kontrprzykład w rozdziale 15}
podaje dokładnie taki funkcjonał na $\ell^2$.
\end{przyklad}

\subsection{Zbieżność słaba i metoda bezpośrednia}

\begin{definicja}[Półciągłość z dołu i koercywność]
\label{def:weak-lsc16}
Niech $H$ będzie rzeczywistą przestrzenią Hilberta,
$C\subseteq H$ niepustym zbiorem i $J:C\to(-\infty,+\infty]$.
Funkcjonał $J$ jest \emph{słabo półciągły z dołu}, jeśli dla każdej sieci
$(x_\alpha)$ w $C$ zbieżnej do $x\in C$ w topologii
słabej $\sigma(H,H')$ zachodzi
$J(x)\le\liminf_\alpha J(x_\alpha)$, gdzie
$\liminf_\alpha J(x_\alpha)=
\sup_\alpha\inf_{\beta\ge\alpha}J(x_\beta)$.
Równoważnie, każdy zbiór podpoziomowy
$\{x\in C:J(x)\le a\}$, $a\in\R$, jest domknięty
w topologii słabej indukowanej na $C$.
Istotnie, warunek dla sieci zapewnia, że granica sieci
z podpoziomu nadal należy do tego podpoziomu.
Odwrotnie, gdyby $\liminf_\alpha J(x_\alpha)<J(x)$,
można byłoby wybrać rzeczywiste $a$ ściśle pomiędzy
tymi wartościami. Otwarte dopełnienie podpoziomu
zawiera $x$, więc od pewnego indeksu $J(x_\alpha)>a$,
co wymusza $\liminf_\alpha J(x_\alpha)\ge a$,
sprzecznie z wyborem $a$.
Analogicznie definiuje się półciągłość z dołu
w topologii normowej.
Funkcjonał jest \emph{koercywny na $C$}, jeśli dla
każdego $a\in\R$ istnieje $R>0$ takie, że
$J(x)>a$ dla każdego $x\in C$ o $\|x\|>R$.
Równoważnie, wszystkie podpoziomy są ograniczone.
\end{definicja}

\begin{lemat}[Podpoziomowe zbiory wypukłe]
\label{lem:convex-weaklsc16}
Niech $H$ będzie rzeczywistą przestrzenią Hilberta,
a $J:H\to(-\infty,+\infty]$ funkcjonałem wypukłym
i półciągłym z dołu w topologii normowej.
Wtedy $J$ jest słabo półciągły z dołu.
\end{lemat}
\begin{proof}
Dla każdego $a\in\R$ zbiór $C_a=\{x:J(x)\le a\}$
jest normowo domknięty przez półciągłość z dołu
i wypukły przez wypukłość $J$. Z Wniosku~\ref{cor:weak-convex-closure}
o równości domknięć zbioru wypukłego wynika,
że każdy normowo domknięty zbiór wypukły jest
domknięty także w topologii słabej. Zatem wszystkie $C_a$
są słabo domknięte, co daje tezę.
\end{proof}

\begin{lemat}[Słaba zbieżność podciągu w przestrzeni Hilberta]
\label{lem:hilbert-weakseq16}
Niech $H$ będzie rzeczywistą przestrzenią Hilberta,
a $(x_n)$ ciągiem ograniczonym w $H$.
Istnieje podciąg $(x_{n_k})$ i $x\in H$ takie,
że $x_{n_k}\to x$ słabo w $H$.
Nie zakładamy ośrodkowości całej przestrzeni $H$.
\end{lemat}
\begin{proof}
Niech $M=\sup_n\|x_n\|<\infty$ i
$H_0=\overline{\operatorname{lin}\{x_n:n\in\N\}}$.
Przestrzeń $H_0$ jest ośrodkowa, ponieważ skończone
kombinacje liniowe $x_n$ o współczynnikach wymiernych
tworzą w niej przeliczalny zbiór gęsty. Jest również
przestrzenią Hilberta jako domknięta podprzestrzeń $H$.
Jeśli $H_0=\{0\}$, cały ciąg jest zerowy i teza zachodzi.
W przeciwnym razie wybieramy bazę ortonormalną $(e_j)$
skończoną lub przeliczalnie nieskończoną.
Dla każdego $j$ ciąg $\langle x_n,e_j\rangle$ jest ograniczony.
Z twierdzenia Bolzana--Weierstrassa wybieramy podciąg,
na którym pierwsza współrzędna ma granicę. Z niego
wybieramy podciąg dla drugiej współrzędnej i tak dalej.
W przypadku skończonej bazy ostatni podciąg ma
wszystkie wymagane granice. W przypadku nieskończonej
bazy wybieramy $k$-ty wyraz $k$-tego podciągu.
Indeksy tak wybranego ciągu rosną, a jego ogon
od $k=j$ jest podciągiem $j$-tego podciągu, więc
$j$-ta współrzędna jest zbieżna dla każdego $j$.

Na przestrzeni skończonych kombinacji wektorów $e_j$
definiujemy $L(v)=\lim_k\langle x_{n_k},v\rangle$.
Istnienie granicy wynika z jej istnienia na każdym $e_j$,
a $|L(v)|\le M\|v\|$. Funkcjonał $L$ przedłuża się
jednoznacznie na $H_0$ jako ciągły funkcjonał liniowy.
Z twierdzenia Riesza istnieje $x\in H_0$ takie,
że $L(v)=\langle x,v\rangle$ dla $v\in H_0$.
Ponieważ $\|x\|=\|L\|\le M$, dla dowolnego $v\in H_0$
i skończonej kombinacji $w$ wektorów $e_j$ zachodzi
\[
 |\langle x_{n_k}-x,v\rangle|
 \le 2M\|v-w\|+|\langle x_{n_k}-x,w\rangle|.
\]
Dla $\varepsilon>0$ wybieramy $w$ tak, aby pierwszy
składnik był mniejszy od $\varepsilon/2$, a następnie
$k$ dostatecznie duże, aby drugi składnik był mniejszy
od $\varepsilon/2$. Gdy $M=0$, ciąg jest zerowy.
Otrzymujemy zbieżność $\langle x_{n_k},v\rangle
\to\langle x,v\rangle$ dla każdego $v\in H_0$.
Dla $v\in H_0^\perp$ obie strony są zerowe.
Dla dowolnego $v\in H$ stosujemy rozkład
$H=H_0\oplus H_0^\perp$ z Twierdzenia~\ref{thm:orth-decomposition13}.
Z twierdzenia Riesza każdy ciągły funkcjonał liniowy
na $H$ jest dany przez iloczyn skalarny z takim $v$.
Otrzymana zbieżność jest więc zbieżnością słabą w $H$.
\end{proof}

\begin{twierdzenie}[Metoda bezpośrednia w przestrzeni Hilberta]
\label{thm:direct-method16}
Niech $H$ będzie rzeczywistą przestrzenią Hilberta,
$C\subseteq H$ niepustym, normowo domkniętym i wypukłym
zbiorem, a $J:C\to(-\infty,+\infty]$ funkcjonałem,
który jest właściwy (tzn. przyjmuje przynajmniej jedną
wartość skończoną), ograniczony z dołu, koercywny na $C$
i półciągły z dołu dla topologii słabej indukowanej na $C$.
Wtedy istnieje $u\in C$ takie, że
$J(u)=\inf_{v\in C}J(v)$.
Jeżeli dodatkowo $J$ jest ściśle wypukły,
minimizator jest jedyny.
\end{twierdzenie}
\begin{proof}
Niech $m=\inf_C J$, który jest skończony na mocy
właściwości i ograniczoności z dołu. Dla każdego $n\ge1$ z definicji infimum wybieramy
$u_n\in C$ takie, że $m\le J(u_n)<m+1/n$.
Wszystkie $u_n$ należą do podpoziomu $\{v\in C:J(v)\le m+1\}$,
który jest ograniczony z koercywności. Zatem
$(u_n)$ jest ograniczonym ciągiem w $H$.
Na mocy Lematu~\ref{lem:hilbert-weakseq16}
pewien podciąg $u_{n_k}\to u$ słabo w $H$.
Ponieważ $C$ jest normowo domknięty i wypukły,
jest również słabo domknięty na mocy
Wniosku~\ref{cor:weak-convex-closure}, więc $u\in C$.
Słaba półciągłość z dołu daje
$J(u)\le\liminf_k J(u_{n_k})=m$; z definicji infimum
zachodzi też $J(u)\ge m$. Jeśli $u\ne v$ są minimizatorami,
ścisła wypukłość daje $J((u+v)/2)<m$, sprzeczność.
\end{proof}

\begin{uwaga}[Czego nie daje sama koercywność?]
\label{rem:coercive16}
W nieskończenie wymiarowej przestrzeni Hilberta
domknięta kula jednostkowa nie jest zwarta w normie:
zawiera ciąg ortonormalny, którego różne wyrazy
są odległe o $\sqrt2$, więc żaden jego podciąg
nie jest zbieżny w normie.
Sama ciągłość normowa i koercywność funkcjonału
nie dają zatem zwartości podpoziomu potrzebnej
do bezpośredniego zastosowania twierdzenia Weierstrassa. W twierdzeniu wykorzystaliśmy
słabą zbieżność podciągów i słabą półciągłość z dołu.
W ogólnej przestrzeni Banacha naturalną rolę przestrzeni
Hilberta przejmują przestrzenie refleksywne.
W rozdziale 26 udowodnimy ogólną wersję twierdzenia
o słabo zbieżnym podciągu ciągu ograniczonego.
Powyższy elementarny dowód w przestrzeni Hilberta
nie wymaga tej późniejszej teorii.
\end{uwaga}

\subsection{Subgradienty i rzutowanie proksymalne}

\begin{definicja}[Subgradient i subróżniczka]
\label{def:subgrad16}
Niech $J:H\to(-\infty,+\infty]$ będzie właściwym
funkcjonałem wypukłym, a $u\in H$ punktem takim,
że $J(u)<+\infty$. Wektor $p\in H$ nazywamy
\emph{subgradientem} $J$ w $u$, jeśli
$J(v)\ge J(u)+\langle p,v-u\rangle$ dla każdego $v\in H$.
Zbiór wszystkich subgradientów oznaczamy $\partial J(u)$
i nazywamy \emph{subróżniczką} $J$ w $u$.
Dla $J(u)=+\infty$ przyjmujemy $\partial J(u)=\varnothing$.
\end{definicja}

\begin{stwierdzenie}[Warunek optymalności i monotoniczność]
\label{prop:subgrad16}
Jeżeli $J:H\to(-\infty,+\infty]$ jest właściwy i wypukły,
to $u$ minimalizuje $J$ wtedy i tylko wtedy, gdy
$0\in\partial J(u)$. Jeśli $p\in\partial J(u)$ oraz
$q\in\partial J(v)$, to
$\langle p-q,u-v\rangle\ge0$.
Jeżeli $J$ jest różniczkowalny Frécheta w $u$ i skończony
na otoczeniu $u$, to $\partial J(u)=\{\nabla J(u)\}$.
\end{stwierdzenie}
\begin{proof}
Pierwsza równoważność wynika wprost z definicji.
Z nierówności subgradientowych
$J(v)\ge J(u)+\langle p,v-u\rangle$ i
$J(u)\ge J(v)+\langle q,u-v\rangle$,
po dodaniu dostajemy monotoniczność.
Jeśli $J$ jest różniczkowalny, to dla $p\in\partial J(u)$
podstawiamy $v=u+th$ z dodatnim i ujemnym $t$,
po czym przechodzimy do zera. Otrzymujemy
$\langle p,h\rangle=DJ(u)h$ dla każdego $h\in H$.
Odwrotnie, dla $v$ z dziedziny efektywnej $J$ i $0<t<1$
wypukłość daje
$(J(u+t(v-u))-J(u))/t\le J(v)-J(u)$.
Przechodząc z $t$ do zera, otrzymujemy
$DJ(u)(v-u)\le J(v)-J(u)$.
Dla $J(v)=+\infty$ nierówność zachodzi automatycznie.
Reprezentacja Riesza daje zatem $\nabla J(u)\in\partial J(u)$.
\end{proof}

\begin{przyklad}[Subróżniczka normy w Hilbercie]
\label{ex:norm-subgrad16}
Dla $J(x)=\|x\|$ zachodzi
$\partial J(0)=\{p:\|p\|\le1\}$ i
$\partial J(x)=\{x/\|x\|\}$ dla $x\ne0$.
Jeżeli $\|p\|\le1$, to $\langle p,v\rangle\le\|v\|$
dla każdego $v\in H$, zatem $p\in\partial J(0)$.
Odwrotnie, dla $p\in\partial J(0)$, podstawiając $v=p$,
otrzymujemy $\|p\|^2\le\|p\|$, czyli $\|p\|\le1$.
Dla $x\ne0$ reguła łańcuchowa zastosowana do
$\|x\|=\sqrt{\langle x,x\rangle}$ daje
$DJ(x)h=\langle x,h\rangle/\|x\|$.
Stwierdzenie~\ref{prop:subgrad16} daje drugi wzór.
W ogólnej przestrzeni Banacha analogiczny subgradient normy
jest funkcjonałem z przestrzeni dualnej, a jego istnienie
dla $x\ne0$ zapewnia twierdzenie Hahna--Banacha.
\end{przyklad}

\begin{twierdzenie}[Odwzorowanie proksymalne]
\label{thm:prox16}
Niech $J:H\to(-\infty,+\infty]$ będzie właściwy,
wypukły, normowo półciągły z dołu i ograniczony z dołu,
a $\lambda>0$. Dla każdego $y\in H$ istnieje dokładnie
jeden punkt $u=\operatorname{prox}_{\lambda J}(y)$
minimalizujący funkcjonał
$v\mapsto J(v)+\|v-y\|^2/(2\lambda)$.
Jest on scharakteryzowany równoważnością
$y-u\in\lambda\partial J(u)$.
Ponadto odwzorowanie
$\operatorname{prox}_{\lambda J}:H\to H$ jest
nieekspansywne:
$\|\operatorname{prox}_{\lambda J}(y)-
\operatorname{prox}_{\lambda J}(z)\|\le\|y-z\|$.
\end{twierdzenie}
\begin{proof}
Funkcjonał $\Phi_y(v)=J(v)+\|v-y\|^2/(2\lambda)$
jest właściwy, wypukły i normowo półciągły z dołu,
a więc słabo półciągły z dołu na mocy
Lematu~\ref{lem:convex-weaklsc16}.
Jest koercywny, gdyż $J$ jest ograniczony z dołu,
i jest ściśle wypukły dzięki składnikowi kwadratowemu.
Twierdzenie~\ref{thm:direct-method16} dla $C=H$
daje jednoznaczny minimizator $u$.

Jeśli $u$ minimalizuje $\Phi_y$, to dla $v$ z
$J(v)<+\infty$ i $0<t<1$, z wypukłości $J$
oraz minimalności mamy
\[
0\le\Phi_y(u+t(v-u))-\Phi_y(u)
\le t[J(v)-J(u)]+\frac t\lambda\langle u-y,v-u\rangle
+\frac{t^2}{2\lambda}\|v-u\|^2.
\]
Po podzieleniu przez $t$ i przejściu z $t$ do zera
otrzymujemy $J(v)\ge J(u)+\langle(y-u)/\lambda,v-u\rangle$,
a więc $(y-u)/\lambda\in\partial J(u)$.
Odwrotnie, ta nierówność oraz rozwinięcie
$\|v-y\|^2$ wokół $u$ dają
$\Phi_y(v)-\Phi_y(u)\ge\|v-u\|^2/(2\lambda)$,
co zapewnia minimalność.

Niech $u=\operatorname{prox}_{\lambda J}(y)$ oraz
$v=\operatorname{prox}_{\lambda J}(z)$.
Monotoniczność subróżniczki daje
$\langle (y-u)-(z-v),u-v\rangle\ge0$.
Stąd $\|u-v\|^2\le\langle y-z,u-v\rangle
\le\|y-z\|\,\|u-v\|$,
co po podzieleniu przez $\|u-v\|$ (gdy $u\ne v$)
daje tezę.
\end{proof}

\subsection{Przykład: słabe zagadnienie brzegowe}

Niech $\lambda$ oznacza miarę Lebesgue'a na $[0,1]$.
Wprowadzamy przestrzeń $H_0^1(0,1)$ przez opis elementarny:
\[
H_0^1(0,1)=\Bigl\{u:[0,1]\to\R:\ 
 u(x)=\int_0^x v(t)\,\mathrm dt\ (0\le x\le1),\ 
 v\in L^2(0,1),\ \int_0^1 v(t)\,\mathrm dt=0\Bigr\}.
\]
Funkcje $u$ są absolutnie ciągłe, ich pochodna $u'=v$
istnieje prawie wszędzie, a warunki brzegowe
$u(0)=u(1)=0$ są spełnione. Przyjmujemy
$\langle u,w\rangle_H=\int_0^1u'(t)w'(t)\,\mathrm dt$.
Odwzorowanie $u\mapsto u'$ jest liniową izometrią
$H_0^1(0,1)$ na domkniętą podprzestrzeń
$\{v\in L^2(0,1):\int_0^1v(t)\,\mathrm dt=0\}$,
przy czym podprzestrzeń ta jest jądrem ciągłego funkcjonału
$v\mapsto\int_0^1v(t)\,\mathrm dt$ (jego norma jest co najwyżej $1$
na mocy nierówności Cauchy'ego--Schwarza).
Surjektywność izometrii wynika z definicji $H_0^1(0,1)$,
a iniektywność z warunku $u(0)=0$ i wzoru całkowego.
Jako domknięta podprzestrzeń przestrzeni Hilberta $L^2(0,1)$
jej obraz jest zupełny, więc $H_0^1(0,1)$ jest przestrzenią Hilberta.

\begin{lemat}[Nierówność Poincarégo na odcinku]
\label{lem:poincare16}
Dla każdego $u\in H_0^1(0,1)$ zachodzi
$\|u\|_{L^2}\le2^{-1/2}\|u'\|_{L^2}$.
\end{lemat}
\begin{proof}
Ponieważ $u(x)=\int_0^xu'(t)\,\mathrm dt$,
nierówność Cauchy'ego--Schwarza daje
$|u(x)|^2\le x\int_0^x|u'(t)|^2\,\mathrm dt
\le x\|u'\|_{L^2}^2$.
Całkując po $x\in[0,1]$, otrzymujemy
$\|u\|_{L^2}^2\le\frac12\|u'\|_{L^2}^2$.
\end{proof}

\begin{twierdzenie}[Zasada minimalizacji dla problemu Dirichleta]
\label{thm:dirichlet16}
Niech $f\in L^2(0,1;\R)$. Na $H=H_0^1(0,1)$
określmy funkcjonał
\[
J(u)=\frac12\int_0^1|u'(t)|^2\,\mathrm dt
-\int_0^1 f(t)u(t)\,\mathrm dt.
\]
Istnieje dokładnie jeden minimizator $u\in H$ funkcjonału $J$.
Jest on dokładnie jedynym elementem $H$ spełniającym
równanie słabe
\[
\int_0^1u'(t)v'(t)\,\mathrm dt
=\int_0^1f(t)v(t)\,\mathrm dt\qquad(v\in H).
\]
Dla $f(t)=1$ rozwiązaniem jest $u(t)=t(1-t)/2$.
\end{twierdzenie}
\begin{proof}
Funkcjonał $J$ jest ciągły i ściśle wypukły na $H$.
Z Lematu~\ref{lem:poincare16} i nierówności
Cauchy'ego--Schwarza wynika
$J(u)\ge\frac12\|u\|_H^2-
2^{-1/2}\|f\|_{L^2}\|u\|_H$, więc jest koercywny
i ograniczony z dołu. Ponieważ jest wypukły
i normowo ciągły, jest słabo półciągły z dołu.
Twierdzenie~\ref{thm:direct-method16} daje minimizator,
a ścisła wypukłość zapewnia jego jednoznaczność.

Dla $h\in H$ rozwinięcie kwadratu daje
\[
J(u+h)-J(u)=\langle u,h\rangle_H-\int_0^1 f(t)h(t)\,\mathrm dt
+\tfrac12\|h\|_H^2.
\]
Część liniowa jest ograniczona przez
$(\|u\|_H+2^{-1/2}\|f\|_{L^2})\|h\|_H$,
a reszta podzielona przez $\|h\|_H$ dąży do zera.
Zatem $J$ jest różniczkowalny Frécheta, a jego pochodna ma postać
$DJ(u)v=\int_0^1u'(t)v'(t)\,\mathrm dt-
\int_0^1 f(t)v(t)\,\mathrm dt$.
Twierdzenie~\ref{thm:convex-derivative16} daje
równoważność $DJ(u)=0$ i globalnej minimalności.
Jeśli $f\equiv1$, funkcja $u(t)=t(1-t)/2$ spełnia
$u(0)=u(1)=0$ oraz $u''(t)=-1$; całkowanie przez części
dla absolutnie ciągłego $v\in H$ daje
$\int_0^1u'v'\,\mathrm dt=\int_0^1v\,\mathrm dt$.
Przez jednoznaczność jest to szukane rozwiązanie.
\end{proof}

\subsection{Wypukłe maksima i twierdzenie Kreina--Milmana}

Twierdzenie Kreina--Milmana i zasada maksimum Bauera orzekają,
że niepusty zwarty wypukły podzbiór przestrzeni lokalnie wypukłej
Hausdorffa jest domkniętą otoczką wypukłą swoich punktów
ekstremalnych, a każda wypukła i półciągła z góry funkcja
rzeczywista na tym zbiorze osiąga maksimum w punkcie ekstremalnym. Ich rola w optymalizacji jest odmienna od metody bezpośredniej:
\emph{zwartość wypukłego zbioru} daje możliwość badania maksimum
funkcjonału wypukłego na punktach ekstremalnych,
natomiast \emph{koercywność i słaba półciągłość z dołu}
służą do znajdowania minimum.

\begin{stwierdzenie}[Maksimum funkcjonału afinicznego]
\label{prop:affine-extreme16}
Niech $E$ będzie lokalnie wypukłą TVS Hausdorffa nad $\K$,
$K\subseteq E$ niepustym, zwartym i wypukłym zbiorem,
a $A:K\to\R$ ciągłą funkcją afiniczną,
czyli $A((1-t)x+ty)=(1-t)A(x)+tA(y)$.
Wtedy istnieje punkt ekstremalny $e\in K$, w którym
$A$ osiąga maksimum, oraz (być może inny) punkt ekstremalny,
w którym osiąga minimum.
\end{stwierdzenie}
\begin{proof}
Funkcje $A$ i $-A$ są wypukłe i ciągłe,
więc zastosowanie zasady maksimum Bauera do obu daje tezę.
\end{proof}

\begin{przyklad}[Miary prawdopodobieństwa i ich punkty ekstremalne]
\label{ex:prob-measures16}
Niech $K$ będzie niepustą zwartą przestrzenią metryczną,
a $\mathcal P(K)$ zbiorem borelowskich miar probabilistycznych na $K$.
Na mocy twierdzenia Riesza--Markowa identyfikujemy
$\mathcal P(K)$ z dodatnimi funkcjonałami na $C(K;\R)$,
przyjmującymi wartość $1$ na funkcji stałej $1$.
W topologii słabej-* $\sigma(C(K)',C(K))$ jest to zbiór
zwarty i wypukły: jako podzbiór kuli jednostkowej
jest domknięty przez warunki
$L(1)=1$ i $L(\varphi)\ge0$ dla każdej $\varphi\ge0$.

Punktami ekstremalnymi $\mathcal P(K)$ są dokładnie
miary Diraca $\delta_x$, $x\in K$.
Dla $\delta_x=t\mu+(1-t)\nu$, $0<t<1$, dodatniość
miar i równość ich mas wymuszają
$\mu(K\setminus\{x\})=\nu(K\setminus\{x\})=0$,
więc $\mu=\nu=\delta_x$.
Jeśli natomiast $\mu$ nie jest miarą Diraca,
istnieje zbiór borelowski $B$ o
$0<\mu(B)<1$: w przeciwnym razie nośnik miary
musiałby być jednopunktowy (dla zwartej metrycznej $K$
wynika to z rozdzielania różnych punktów nośnika
rozłącznymi otoczeniami o dodatniej mierze).
Wówczas
$\mu=\mu(B)\,\mu_B+(1-\mu(B))\,\mu_{K\setminus B}$,
gdzie $\mu_B$ i $\mu_{K\setminus B}$ są znormalizowanymi
ograniczeniami $\mu$, a obie miary są różne.
Twierdzenie Kreina--Milmana stwierdza więc, że
$\mathcal P(K)$ jest słabo-* domkniętą otoczką wypukłą
rodziny $\{\delta_x:x\in K\}$.
W szczególności dla $\varphi\in C(K;\R)$
maksimum funkcjonału $\mu\mapsto\int_K\varphi(x)\,\mathrm d\mu(x)$
można osiągnąć na pewnej mierze Diraca.
\end{przyklad}

\begin{uwaga}[Dalsza perspektywa]
Twierdzenie Kreina--Milmana jest rezultatem o \emph{zwartych
zbiorach wypukłych} w ogólnych przestrzeniach lokalnie wypukłych,
a nie jedynie o kulach jednostkowych przestrzeni Banacha.
Z kolei subróżniczki i operatory proksymalne prowadzą
do teorii operatorów monotonicznych i metod minimalizacji,
co wykracza poza podstawowy zakres obecnego rozdziału.
\end{uwaga}
